%% -------------------------------------------------------------------
\subsection{\vrev{BCM 운영 통합 프레임워크}}
\label{subsec:bcm-integration}
%% -------------------------------------------------------------------

2단계 프레임워크와 에스컬레이션 체계를 기존 BCMS에 통합하기 위한
구체적 운영 절차를 다음과 같이 제시한다.

\subsubsection{\tmp{데이터 수집 체계}}

\inew{$\DRTO$ 모델의 두 입력 변수인 \textbf{관측성 $O(t)$} 와 \textbf{불확실성 $U(t)$} 의
수집은 각각 기존 \chg{운영과 보안} 모니터링 인프라와 \chg{장애 이력 및 위협 인텔리전스} 데이터의
연동을 통해 구현한다.}

\inew{\chg{관측성 지표 $O(t)$의 수집은 기존 모니터링 인프라와의 연동을 통해 다음과 같이 구현한다. 첫째, APM(Application Performance Monitoring) 연동으로 서비스 응답 시간, 오류율, 처리량을 종합하여 $O_{\mathrm{APM}} \in [0,1]$로 정규화하며, 가중 평균은 $O_{\mathrm{APM}} = w_1 \cdot R_{\mathrm{avail}} + w_2 \cdot (1-R_{\mathrm{error}}) + w_3 \cdot R_{\mathrm{perf}}$($\sum w_i = 1$)로 산출한다. 둘째, SIEM(Security Information and Event Management) 연동으로 보안 이벤트 심각도와 대응 상태를 기반으로 $O_{\mathrm{SIEM}} \in [0,1]$을 산출하되, 미대응 고위험 이벤트가 존재하면 감점한다. 셋째, 종합 관측성은 $O(t) = \alpha \cdot O_{\mathrm{APM}} + (1-\alpha) \cdot O_{\mathrm{SIEM}}$로 정의하며, $\alpha$는 조직 특성에 따라 예를 들어 $[0.5, 0.8]$에서 설정한다.}}

\inew{\chg{불확실성 지표 $U(t)$의 수집은 조직 내부의 과거 장애 이력과 외부 위협 인텔리전스 데이터의 연동을 통해 다음과 같이 구현한다. 첫째, 장애 이력 DB 연동으로 조직 자체의 과거 장애와 복구 이력(MTTR, RTO 초과 사례, 복구 소요 시간 분포)을 적합 분포(가우시안 또는 파레토)로 모델링하여 $U_{\mathrm{hist}} \in [0,\; U_{\max}]$을 산출하며, 분포 형태 판정에는 Hill 추정량~\citep{hill1975simple} 또는 Anderson-Darling 검정~\citep{anderson1954test}을 적용하여 C2(경꼬리) 또는 C3(중꼬리) 환경으로 자동 분기한다. 둘째, 위협 인텔리전스(CVE/MITRE ATT\&CK) 연동으로 자산 대장과 매핑된 CVE 공개 DB, EPSS 점수, MITRE ATT\&CK 위협 액터 활동 정보를 기반으로 미래 위험 강도 $U_{\mathrm{threat}} \in [0,\; U_{\max}]$을 산출하되, 미패치 고위험 CVE 수와 EPSS 가중 평균을 결합한다. 셋째, 종합 불확실성은 $U(t) = \beta \cdot U_{\mathrm{hist}} + (1-\beta) \cdot U_{\mathrm{threat}}$로 정의하며, $\beta$는 조직의 과거 데이터 풍부도와 위협 환경 변동성에 따라 $[0.5, 0.8]$에서 설정한다. 장애 이력이 풍부한 성숙 조직은 $\beta$가 높고, 신생 조직과 신규 위협 환경은 $\beta$가 낮다.}}

\inew{원시 변수 $O_{\mathrm{raw}}$, $U_{\mathrm{raw}}$, 가중치 $k_O$의 제1계층
하위 지표 조작화 예시와 가상 기업 라이프사이클 적용 시나리오는
\secref{app:operationalization}(관측성 $O_{\mathrm{raw}}$ 조작화 예시,
불확실성 $U_{\mathrm{raw}}$ 조작화 예시)에 참고용으로 수록되어 있다.}

\subsubsection{\tmp{$\eta$ 모니터링 대시보드}}

$\eta$ 비율의 실시간 시각화를 통해 BCMS 효과성을 지속적으로 측정한다.
대시보드의 핵심 구성요소는 \chg{실시간 $\eta$ 추이 차트(시계열),
에스컬레이션 등급 표시(Green/Yellow/Orange/Red),
$O(t)$와 $U(t)$의 개별 추이, CVaR$_{99}$ 현재값}이며,
이는 \nrev{ISO 22301:2019} 조항~9.1(모니터링, 측정, 분석, 평가)의
정량적 요구사항을 직접 충족한다.

\subsubsection{\tmp{사후 분석 프로세스}}

복구 완료 후 24--72시간 이내에 다음의 사후 분석을 수행한다.
\chg{먼저 실제 복구 시간 $T_{\mathrm{actual}}$과 $\DRTO$ 예측값을 비교하고,
절대 오차 $E = T_{\mathrm{actual}} - \DRTO_{\mathrm{final}}$과
상대 오차 $E_r = E / T_{\mathrm{actual}}$를 산출하며,
과대 및 과소 추정 원인을 분석하여 파라미터 교정을 권고하고,
장애 이력 DB 갱신 및 $U(t)$ 분포 재추정을 수행한다.}
이 프로세스는 \nrev{ISO 22301:2019} 조항~10.1(지속적 개선)과 직접 대응된다.


